\chapter{Persistenza dei dati: JPA}

\section{Interfacciarsi a un DB relazionale}

Come si eseguono query SQL da un'applicazione? Un DB Server espone i suoi servizi tramite:
una \emph{CLI} (command line interface), eventuali \emph{GUI client}, e soprattutto una
\textbf{API} che rende i servizi accessibili ad applicazioni esterne (locali o remote via
TCP). L'API \`e \emph{specifica del DB} e viene resa disponibile tramite \textbf{driver}:
librerie che rendono le operazioni invocabili in un dato linguaggio (driver C++, Python,
JavaScript, Java\dots).

\begin{definizione}[JDBC]
\textbf{JDBC} \`e un insieme di classi e interfacce a cui i driver Java \emph{devono
conformarsi}, rendendo il codice Java \textbf{indipendente dal DB}: baster\`a installare il
driver giusto (per Postgres, MySQL, MariaDB\dots) senza cambiare il codice.
\end{definizione}

\section{ORM, JPA, Hibernate}

\begin{definizione}[ORM]
L'\textbf{ORM} (Object-Relational Mapping) \`e una tecnica per mappare i costrutti dei DB
relazionali in costrutti object-oriented. Uno strumento di ORM fornisce un'API che permette
di accedere alla persistenza \emph{con un modello a oggetti} anzich\'e relazionale: si legge
e scrive creando/eliminando oggetti e invocando metodi, anzich\'e con query SQL.
\end{definizione}

\begin{definizione}[JPA e Hibernate]
\textbf{JPA} (Jakarta Persistence API) \`e parte di Jakarta EE: una \emph{specifica} che
descrive come gestire dati relazionali da Java. JPA \textbf{non implementa} un ORM: \`e la
specifica dell'interfaccia che un ORM per Java deve rispettare -- cos\`i lo sviluppatore usa
l'API di JPA senza essere vincolato a uno specifico \emph{provider}. \textbf{Hibernate} \`e
il provider JPA preferito da Spring Boot (esiste anche come ORM indipendente); la cosa pi\`u
semplice \`e richiedere via Maven il set \texttt{spring-boot-starter-data-jpa}, che include
proprio Hibernate. Hibernate interagisce col DB tramite \textbf{JDBC}: fra le dipendenze va
incluso anche il \emph{driver JDBC} del DBMS scelto (nel corso: Postgres).
\end{definizione}

Lo ``stack'' completo: \emph{Spring App} $\rightarrow$ \emph{JPA} $\rightarrow$
\emph{Hibernate ORM} $\rightarrow$ \emph{JDBC} $\rightarrow$ \emph{DB}.

\section{Le @Entity: descrivere entit\`a e relazioni}

\begin{definizione}[Entity]
Un'\textbf{entit\`a} del dominio in JPA corrisponde a una \emph{classe}, tipicamente mappata
su una \emph{tabella} del DB; le istanze sono mappate su uno o pi\`u record; i campi
esprimono propriet\`a rappresentabili come colonne (o come tabelle relazionali, per i campi
che esprimono relazioni). La classe deve avere un \textbf{costruttore di default}: pu\`o
averne altri per il nostro codice, ma quando JPA crea le istanze caricandole dal DB usa
quello di default e poi imposta i campi. Si descrivono entity e relationship con le
annotazioni; \emph{della traduzione in tabelle si occupa JPA} -- ma conoscere la
rappresentazione sottostante aiuta a sfruttarla al meglio.
\end{definizione}

\subsection{Le annotazioni principali}
\begin{itemize}
  \item \textbf{@Entity}: la classe \`e un'entit\`a del DB gestita da JPA;
  \item \textbf{@Table(name = "table\_name")}: dettagli della tabella (di solito il nome);
  \item \textbf{@Id}: il campo \`e la \emph{chiave primaria};
  \item \textbf{@GeneratedValue(strategy = GenerationType.STRATEGY)}: la chiave \`e generata
        dal DBMS, con la strategia indicata; per semplicit\`a si usa
        \texttt{GenerationType.AUTO}, che demanda la scelta al DBMS;
  \item \textbf{@Column}: dettagli della colonna -- \texttt{name} (se diverso dal campo),
        \texttt{nullable} (true/false), \texttt{unique} (true/false).
\end{itemize}

\subsection{Le relazioni}
\begin{itemize}
  \item \textbf{Uno-a-Molti} fra A (l'uno) e B (i molti): si annota un campo in B con
        \textbf{@ManyToOne} (di tipo A) e un campo simmetrico in A con \textbf{@OneToMany}
        (una \texttt{Collection} di B), specificando \texttt{mappedBy = "owner"} dove
        \emph{owner} \`e il nome del campo di B annotato @ManyToOne;
  \item \textbf{Molti-a-Molti} fra A e B: campi simmetrici in entrambe le entity, entrambi
        \textbf{@ManyToMany}, entrambi Collection (\emph{meglio Set che List}); uno dei due
        \`e l'\emph{owner} e aggiunge \textbf{@JoinColumn} che specifica la \emph{join table}
        da usare; l'altro indica con \texttt{mappedBy} il nome del campo che ``possiede'' la
        relazione.
\end{itemize}

\section{I @Repository e le operazioni CRUD}

\begin{definizione}[Repository]
Per operare su una tabella si definisce un \textbf{@Repository}: un'\emph{interface} che
estende \texttt{JpaRepository<E, Id>} (generico nella classe entity E e nel tipo Id della
chiave primaria). JpaRepository definisce gi\`a i metodi CRUD di base, quindi spesso
l'estensione pu\`o restare \emph{vuota}:
\begin{itemize}
  \item \texttt{E save(E entity)}: salva l'entity e \emph{restituisce l'entity risultante}
        (che pu\`o essere diversa da quella di partenza!). Se l'id \`e gi\`a nel DB \`e una
        \emph{update}, altrimenti una \emph{create};
  \item \texttt{Iterable<E> saveAll(Iterable<E>)}: salva e restituisce un elenco;
  \item \texttt{Iterable<E> findAll()}: tutte le entity della tabella;
  \item \texttt{Optional<E> findById(Id id)}: l'entity con quell'id, se presente;
  \item \texttt{void deleteById(Id id)} / \texttt{void delete(E entity)}: eliminano, se
        presente.
\end{itemize}
\end{definizione}

\subsection{Le query derivate}
Aggiungendo al repository dichiarazioni di metodi che seguono una precisa
\emph{naming convention}, JPA \textbf{deriva automaticamente la query SQL} che li implementa:
\begin{center}
\texttt{findByAuthorAndTitleContainsOrderByTitleAsc(String author, String title)}\\
\texttt{deleteTop1ByAuthorBeforeIgnoreCase(String author)}\\
\texttt{countDistinctByTitleStartsWith(String title)}
\end{center}
Struttura del nome:
\begin{itemize}
  \item \textbf{operazione}: cosa si vuole fare (sostanzialmente \texttt{find},
        \texttt{delete}, \texttt{count});
  \item \textbf{modificatore dell'operazione}: limita o condiziona i risultati
        (\texttt{Top1}, \texttt{Distinct}\dots);
  \item \textbf{By}: separa il ``soggetto'' della query dal predicato;
  \item \textbf{specificatori di predicato}: una colonna + un'operazione logica (se omessa,
        uguaglianza) per filtrare in base al valore; \emph{per ogni predicato serve un
        parametro in ingresso} col valore da verificare;
  \item \textbf{operatori logici}: \texttt{And} / \texttt{Or} combinano pi\`u specificatori;
  \item \textbf{modificatori del predicato}: \texttt{IgnoreCase} (sullo specificatore che lo
        precede), \texttt{AllIgnoreCase} (sull'intero predicato), \texttt{OrderBy} + colonne
        + \texttt{Asc}/\texttt{Desc} (ordinamento dei risultati).
\end{itemize}

\section{Il ciclo di vita delle @Entity}

\begin{importante}[Gli stati: new, managed, detached, removed]
\begin{enumerate}
  \item \textbf{New}: l'entit\`a \`e appena creata in memoria (con \texttt{new}): JPA non la
        gestisce e \emph{il DB non ha idea che esista}; le modifiche ai campi non hanno
        effetto sul DB.
  \item \textbf{Managed}: quando l'entit\`a viene salvata con \texttt{repo.save(e)} (create)
        diventa \emph{gestita}; lo stesso accade \emph{estraendola dal DB} con un metodo del
        repository (es. \texttt{findById}). Finch\'e \`e managed, \textbf{tutte le modifiche
        ai campi dell'oggetto si riflettono automaticamente nel DB} (il cosiddetto
        \emph{dirty checking}). Resta managed \emph{finch\'e non termina la transazione
        corrente}.
  \item \textbf{Detached}: al termine della transazione l'entit\`a \`e ``staccata'': le
        modifiche \emph{non} si riflettono pi\`u sul DB. Per riagganciarla si chiama di nuovo
        \texttt{save}: poich\'e l'id \`e gi\`a nel DB, la save far\`a una \emph{update}
        (anzich\'e una create).
  \item \textbf{Removed}: cancellando con \texttt{repo.delete(e)} l'entit\`a \`e ``rimossa'':
        non appena la transazione termina, scompare dal DB.
\end{enumerate}
\end{importante}

\begin{attenzione}[Quando termina una transazione?]
Se non altrimenti specificato, la transazione \`e \textbf{limitata ai singoli metodi del
repository}: chiamando \texttt{repo.save(e)}, l'entit\`a \`e managed \emph{durante}
l'esecuzione di save, ma appena save restituisce \emph{non lo \`e pi\`u} (subito detached).
Lo stesso per \texttt{findById} e gli altri metodi di ricerca.
\end{attenzione}

\subsection{@Transactional}
\begin{definizione}[@Transactional]
Un metodo \textbf{@Transactional} garantisce che le entit\`a restino managed \emph{fino al
termine del metodo}: le modifiche vengono trasferite nel DB (\emph{commit}) al termine se il
metodo ha successo, o annullate (\emph{rollback}) se fallisce. \`E per questo che i metodi di
un @Service sono \emph{tipicamente} @Transactional.
\end{definizione}
\begin{attenzione}
La transazione \`e effettiva \textbf{solo se il metodo \`e invocato da un altro componente
che inietta il servizio via Dependency Injection}: solo cos\`i il metodo viene eseguito
``sotto il controllo'' di Spring, che pu\`o farlo precedere e seguire dalle operazioni che
garantiscono la transazionalit\`a.
\end{attenzione}

\subsection{save: persist o merge}
\texttt{save} pu\`o fare due cose diverse:
\begin{itemize}
  \item \textbf{persist}, se l'entit\`a \emph{non ha id}: aggiorna direttamente l'oggetto
        fornito in input;
  \item \textbf{merge}, se l'entit\`a \emph{ha gi\`a un id}: recupera dal DB una versione
        managed, la aggiorna coi valori della versione detached, e \emph{restituisce la nuova
        entit\`a recuperata}.
\end{itemize}
\begin{attenzione}
In alcuni casi save restituisce lo \emph{stesso} oggetto ricevuto, in altri \textbf{no}:
buona norma lavorare sempre con l'entit\`a \emph{restituita} da save.
\end{attenzione}
